\documentclass[12pt]{article}

\setlength{\parindent}{0pt}
\setcounter{secnumdepth}{3}
\usepackage[margin=1in]{geometry}

% packages
\usepackage{amsmath, amsfonts, amsthm, amssymb}
\usepackage{graphicx, float}
\usepackage{mathtools}
\usepackage{titlesec}
\usepackage{interval}
\usepackage{hyperref}
\usepackage{siunitx}
\usepackage{titling}
\usepackage{vwcol}
\usepackage{setspace}
\usepackage{empheq}
\usepackage{cancel}
\usepackage{esdiff}
\usepackage{multicol}
\usepackage{mdframed}
\usepackage{esdiff}
\usepackage{tikzsymbols}
\usepackage{multicol}
\usepackage{tikz}
\usepackage{varwidth}
\usepackage{pgfplots}
\usepackage{fancyhdr}
\usepackage{enumitem}
\usepackage{tcolorbox}

\intervalconfig {
	soft open fences
}

% header
\pagestyle{fancy}
\fancyhf{}
\lhead{Ethan Fan}
\rhead{AP Calculus BC}
\cfoot{\thepage}

% section format
\titleformat{\section}
  {\large \bfseries}
  {}
  {0em}
  {}
\titlespacing*{\section}{0pt}{0.5em}{0.3em}

\titleformat{\subsection}[runin]
  {\bfseries}
  {}
  {0em}
  {}
\titlespacing*{\subsection}{0pt}{0pt}{0em}

\titleformat{\subsubsection}[runin]
  {\itshape}
  {(\roman{subsubsection})}
  {0em}
  {}

% commands
\newcommand{\alignedintertext}[1]{%
  \noalign{%
    \vskip\belowdisplayshortskip
    \vtop{\hsize=\linewidth#1\par
    \expandafter}%
    \expandafter\prevdepth\the\prevdepth
  }%
}

\newcommand*{\paren}[1]{\ensuremath\left(#1\right)}
\newcommand*{\limit}[2][x]{\ensuremath{\displaystyle\lim_{#1\to#2}}}
\newcommand*{\Limit}[3][x]{\ensuremath{\displaystyle\lim_{#1\to#2}\left[#3\right]}}
\newcommand*{\deriv}[1][x]{\ensuremath{\dfrac{\mathrm{d}}{\mathrm{d}#1}}}
\newcommand*{\Deriv}[2][x]{\ensuremath{\dfrac{\mathrm{d}}{\mathrm{d}#1}\left[#2\right]}}
\newcommand{\ddx}{\dfrac{d}{dx}}
\newcommand{\dydx}{\dfrac{dy}{dx}}
\newcommand{\dydt}{\dfrac{dy}{dt}}
\newcommand{\dxdt}{\dfrac{dx}{dt}}
\newcommand*{\iinteg}[2][x]{\ensuremath{\displaystyle\int #2\;\mathrm{d}#1}}
\newcommand*{\dinteg}[4][x]{\ensuremath{\displaystyle\int_{#2}^{#3}#4\;\mathrm{d}#1}}
\newcommand*{\abs}[1]{\ensuremath{\left|#1\right|}}
\newcommand*{\eps}{\varepsilon}
\newcommand*{\real}{\ensuremath\mathbb{R}}
\newcommand*{\integer}{\ensuremath\mathbb{Z}}
\newcommand*{\posint}{\ensuremath\mathbb{Z^+}}
\newcommand*{\cbrt}[1]{\ensuremath{\sqrt[3]{#1}}}
\newcommand*{\lheqzero}{\ensuremath{\underset{\text{L'H}}{\overset{\left[\frac00\right]}{=}}}}
\newcommand*{\lheqinfty}{\ensuremath{\underset{\text{L'H}}{\overset{\left[\frac{\infty}{\infty}\right]}{=}}}}
\newcommand{\tor}{\text{ or }}
\newcommand{\floor}[1]{\lfloor #1 \rfloor}
\newcommand{\at}{\Big|}

\newcommand{\vem}{\vspace{0.5em}}
\newcommand{\example}[1]{\section{Example #1}}
\newcommand{\problem}[1]{\section{Problem #1}}
\newcommand{\subproblem}[1]{\subsection{(#1)} \,}

\DeclareMathOperator{\dne}{DNE}
\DeclareMathOperator{\sgn}{sgn}

\newcommand{\definition}[1]{\begin{tcolorbox}[colback=red!5!white,colframe=red!75!black,parbox=false] #1 \end{tcolorbox}}

\newcommand{\theorem}[2]{\begin{tcolorbox}[title={#1},colback=blue!5!white,colframe=blue!75!black,parbox=false] #2 \end{tcolorbox}}

\allowdisplaybreaks
% \postdisplaypenalty=100000

% document
\begin{document}

\vspace*{0cm}

% opening
\begin{center}
    {\Large \bfseries Problem Set 23} \\[0.5cm]
    {\large Volumes IV: Cylindrical Shells} \\[0.35cm]
    {\normalsize November 2, 2026}
\end{center}

% problems

\problem{1}

\subproblem{a}
\begin{gather*}
    V=\pi h (r_2^2-r_1^2) = 2\pi \cdot \frac{r_1+r_2}{2}\cdot h\cdot (r_2-r_1)=2\pi rh\Delta r
\end{gather*}
No, the formula would not work if $r$ were the inner radius. \vem

\subproblem{b}
\begin{gather*}
    V=2\pi \cdot 25\cdot 0.1\cdot 2.05=10.25 \pi \\
    V=10\pi 
\end{gather*}

\subproblem{c}
$2\pi \overline{x}_i$ is the circumference. $f(\overline{x}_i)$ is the height. $\Delta x$ is the width.
\begin{gather*}
    V= \lim_{n\to \infty} \sum_{i=1}^n 2\pi \overline{x}_i\, f(\overline{x}_i)\, \Delta x =\int_a^b 2\pi x \, f(x)\, dx
\end{gather*}

\problem{2}

\subproblem{b}
\begin{gather*}
    y=x-x^2\\
    V=2\pi \int_0^1 x^2-x^3\, dx\\
    =2\pi (\frac{x^3}{3}-\frac{x^4}{4})\Big|_0^1\\
    = \boxed{\frac{\pi}{6}}
\end{gather*}

\subproblem{c}
\begin{gather*}
    V=2\pi \int_0^1 x(x^3-2x^2+x)\, dx\\
    =2\pi (\frac{x^5}{5}-\frac{x^4}{2}+ \frac{x^3}{3})\Big|_0^1\\
    = \boxed{\frac{\pi}{15}}
\end{gather*}

\problem{3}

\subproblem{b}
\begin{gather*}
    x=\cbrt y\\
    V=2\pi \int_0^8 y\cdot\cbrt y\, dy\\
    = 2\pi \cdot \frac{3}{7} y^\frac{7}{3}\Big|_0^8\\
    = \boxed{\frac{768\pi}{7}}\\
    A= 64-x^6\\
    V=\pi \int_0^264- x^6\, dx\\
    = \pi ( 64x-\frac{x^7}{7} )\Big|_0^2\\
    = \boxed{\frac{768\pi}{7}}
\end{gather*}

\problem{4}

\subproblem{b}
\begin{gather*}
    V=2\pi \int_0^1 (2-x) x^4\, dx\\
    =2\pi \int_0^1-x^5+2x^4\, dx\\
    =2\pi (\frac{2}{5}x^5-\frac{x^6}{6})\Big|_0^1\\
    = \boxed{\frac{7\pi}{15}}
\end{gather*}

\subproblem{d}
\begin{gather*}
    x=\cbrt y\\
    V= 2\pi\int_0^1 (1-x)(1-\cbrt x)\, dx\\
    =2\pi\int_0^1 x^\frac{4}{3}-x-x^\frac{1}{3}+1\, dx\\
    =2\pi (\frac{3}{7}x^\frac{7}{3}-\frac{x^2}{2}-\frac{3}{4}x^\frac{4}{3}+x)\Big|_0^1\\
    =\boxed{\frac{5\pi}{14}}
\end{gather*}

\problem{5}

\subproblem{a}
\begin{gather*}
    V=2\pi \int_1^3 x(x^2-4x+7-4(x-2)^2)\, dx\\
    = \begin{bmatrix}
        u = x-2\\
        du=dx
    \end{bmatrix} 2\pi \int_{-1}^1 (u+2)(3-3u^2)\, du\\
    = 2\pi \int_{-1}^1 -3u^3-6u^2+3u+6\, du\\
    = 2\pi (-\frac{3}{4}u^4-2u^3+\frac{3}{2}u^2+6u)\Big|_{-1}^1\\
    =\boxed{16\pi }
\end{gather*}

\subproblem{b}
\begin{gather*}
    V=2\pi \int_0^1 xe^{-x^2}\, dx\\
    = \begin{bmatrix}
        u=x^2\\
        du=2x\, dx
    \end{bmatrix} \pi \int_0^1 e^{-u}\, du\\
    = \pi (e^{-u})\Big|_1^0\\
    = \boxed{\pi (1-\frac{1}{e})}
\end{gather*}
The disk method is a different story as it is simply too hard to integrate.\vem

\problem{6}

\subproblem{a}
\begin{gather*}
    \boxed{V=2\pi \int_1^2 x\ln x \, dx}
\end{gather*}

\subproblem{c}
\[
\boxed{V=2\pi \int_0^\pi (4-y)\sqrt{\sin y}\, dy}
\]

\problem{7}

\subproblem{a}
The solid obtained by rotating the region under $y=x^4$ on $(0,3)$ about the $y$-axis.\vem

\subproblem{b}
The solid obtained by rotating the region bounded by $x=1-y^2$, $x=0$, $y=0$, and $y=1$ about the line $y=3$.

\problem{8}

\subproblem{a}
\begin{gather*}
    x^2+y^2=r^2\implies y=\sqrt{r^2-x^2}\\
    V=4\pi \int_0^r x\sqrt{r^2-x^2}\, dx\\
    = \begin{bmatrix}
        u=r^2-x^2\\
        du=-2x\, dx
    \end{bmatrix} -2\pi \int_{r^2}^0 \sqrt{u}\, du\\
    = 2\pi \cdot \frac{2}{3}u^\frac{3}{2} \Big|_0^{r^2}\\
    =\boxed{\frac{4}{3}\pi r^3}
\end{gather*}

\subproblem{b}
\begin{gather*}
    y=-\frac{h}{r}x+h\\
    V=2\pi \int_0^r x(-\frac{h}{r}x+h)\, dx\\
    = 2\pi \int_0^r -\frac{h}{r}x^2+hx\, dx\\
    =2\pi (-\frac{h}{3r}x^3+\frac{h}{2}x^2)\Big|_0^r\\
    =2\pi(\frac{1}{2}r^2h - \frac{1}{3}r^2h)\\
    = \boxed{\frac{1}{3}\pi r^2h}
\end{gather*}

\subproblem{c}
\begin{gather*}
    x^2+y^2=r^2\implies x=\sqrt{r^2-\frac{H^2}{4}}\\
    V=4\pi \int_{\sqrt{r^2-\frac{H^2}{4}}}^rx\sqrt{r^2-x^2} \, dx\\
     = \begin{bmatrix}
        u=r^2-x^2\\
        du=-2x\, dx
    \end{bmatrix} 2\pi \int_{0}^{\frac{H^2}{4}} \sqrt{u}\, du\\
    = 2\pi \cdot \frac{2}{3}u^\frac{3}{2}\Big|_{0}^{\frac{H^2}{4}} \\
    = \boxed{\frac{\pi}{6} H^3}
\end{gather*}

\problem{9}

\subproblem{a}
\begin{gather*}
    A=\int_0^\pi \sin x \, dx = \boxed{2}    
\end{gather*}

\subproblem{b}
\begin{gather*}
    V=2\pi \int_0^\pi x\sin x\, dx\\
    =\begin{bmatrix}
        u=x & dv=\sin x\, dx\\
        du=dx & v=-\cos x
    \end{bmatrix} 2\pi x\cos x \Big|_\pi ^0 +2\pi \int_0^\pi \cos x\, dx\\
    =2\pi^2 + 2\pi \cdot 0= \boxed{2\pi^2}
\end{gather*}

\subproblem{c}
\begin{gather*}
    V=\pi \int_0^\pi \sin ^2 x\, dx\\
    = \frac{\pi}{2} \int_0^\pi  1-\cos  2x\, dx\\
    =\frac{\pi }{2}(x-\frac{\sin 2x}{2})\Big|_0^\pi\\
    = \boxed{\frac{\pi^2}{2}}
\end{gather*}

\subproblem{d}
\begin{gather*}
    V=\boxed{2\pi^2}
\end{gather*}
The region is symmetric about $x=\frac{\pi}{2}$. \vem

\subproblem{e}
\begin{gather*}
    V=\pi \int_0^\pi (\sin x+1)^2-1\, dx\\
    =\pi \int_0^\pi \sin^2 x+2\sin x\,d x\\
    =\pi  \int_0^\pi  \frac{1}{2}-\frac{1}{2}\cos 2x+2\sin x \, dx\\
    =\pi (\frac{x}{2}-\frac{\sin 2x}{4}-2\cos x)\at _0^\pi\\
    =\boxed{\frac{\pi^2}{2}+4\pi}
\end{gather*}

\subproblem{f}
\begin{gather*}
    V=A\cdot 2\pi \overline{x} \implies \overline{x} = \frac{\pi}{2}\\
    V=A\cdot 2\pi \overline{y} \implies \overline{y}= \frac{\pi}{8}\\
    \boxed{\paren{\frac{\pi}{2},\frac{\pi}{8}}}
\end{gather*}

\problem{10}

\subproblem{a}
\begin{gather*}
    y=f(x)\implies x=g(y)\\
    \boxed{V=\int_0^{f(b)}\pi (b^2-g(y)^2)\, dy}
\end{gather*}

\subproblem{b}
\begin{gather*}
    y=f(x)\implies dy=f'(x) \, dx\\
    V=\int_0^{f(b)}\pi (b^2-g(f(x))^2)\, dy\\
    = \pi b^2 f(b)- \int_0^{b} \pi x^2\, dy\\
    =\pi b^2 f(b)- \int_0^{b} \pi x^2f'(x)\, dx\\
    \int_0^{b} \pi x^2f'(x)\, dx\\
    =\begin{bmatrix}
        u=x^2 & dv=f'(x)\, dx\\
        du=2x\, dx  & v=f(x)
    \end{bmatrix} \pi x^2f(x)\at_0^b - 2\pi \int_0^b x f(x) \, dx\\
    \boxed{V= \int_0^b 2\pi xf(x)\, dx}
\end{gather*}

\subproblem{c}
$\pi b^2f(b)$ represents the volume of a cylinder with radius $b$ and height $f(b)$. $\int_0^b \pi x^2 f'(x)\, dx$ represents the volume of the solid formed by rotating the region bounded by $y=f(x)$, $x=0$, and $y=f(b)$ around the $y$-axis.\vem

\subproblem{d}
\begin{gather*}
    V=2\pi \overline{x} A\\
    =2\pi \cdot \frac{1}{A}\int_a^bx f(x)\, dx\cdot A\\
    =\boxed{ \int_a^b2\pi xf(x)\, dx}
\end{gather*}

\subproblem{e}
\begin{gather*}
    f(0)=0 \implies y=0\\
    f(0)=c>0:\\
    V=\pi b^2 c+\int_c^{f(b)}\pi (b^2-g(y)^2)\, dy
\end{gather*}















\end{document}

